% =============================================================================
% Membrane Cosmology v4.7.8 (English main body)
%   Galactic Rotation Without Dark Matter from Elastic Membrane Cosmology:
%   The Geometric Mean Law and Observational Tests
%
%   Sakaguchi Shinobu, 2026-04
%   Rebuild from ReportLab PDF (57 KB, 13 pages, English)
%   Target: arXiv astro-ph.CO
%
%   Revision log:
%     v4.7.8 (2026-04, original): "rev. LV independence"
%     v4.7.8 (2026-04-21, this): affiliation correction
%       Kobe, Japan -> Shiso, Hyogo, Japan (correct address: 兵庫県宍粟市)
%
%   Compile: pdflatex membrane_v478 && bibtex membrane_v478 && pdflatex x2
% =============================================================================
\documentclass[11pt,a4paper]{article}

% ---------------- core packages ---------------------------------------------
\usepackage[a4paper,margin=22mm,headheight=14pt]{geometry}
\usepackage[T1]{fontenc}
\usepackage[utf8]{inputenc}
\usepackage{lmodern}
\usepackage{textcomp}
\usepackage{amsmath,amssymb,amsthm}
\usepackage{mathtools}
\usepackage{booktabs}
\usepackage{array}
\usepackage{tabularx}
\usepackage{longtable}
\usepackage{graphicx}
\usepackage{xcolor}
\usepackage{titlesec}
\usepackage{enumitem}
\usepackage[hidelinks,breaklinks=true]{hyperref}
\usepackage[round,authoryear,sort&compress]{natbib}

% ---------------- section formatting ----------------------------------------
\titleformat{\section}{\normalfont\Large\bfseries}{\thesection.}{0.6em}{}
\titleformat{\subsection}{\normalfont\large\bfseries}{\thesubsection}{0.6em}{}
\titleformat{\subsubsection}{\normalfont\normalsize\bfseries}{\thesubsubsection}{0.6em}{}
\titlespacing*{\section}{0pt}{3.2ex plus 1ex minus .2ex}{1.6ex plus .2ex}
\titlespacing*{\subsection}{0pt}{2.4ex plus .8ex minus .2ex}{1.2ex plus .2ex}

% ---------------- typography tuning ----------------------------------------
\setlength{\emergencystretch}{3em}  % handle occasional overfull hbox
\tolerance=1000
\hbadness=5000

% ---------------- number formatting macros ----------------------------------
\newcommand{\e}[1]{\times 10^{#1}}

% ---------------- physics symbol shortcuts ----------------------------------
\newcommand{\gobs}{g_{\rm obs}}
\newcommand{\gbar}{g_{\rm bar}}
\newcommand{\gN}{g_{N}}
\newcommand{\gc}{g_{c}}
\newcommand{\ao}{a_{0}}
\newcommand{\vflat}{v_{\rm flat}}
\newcommand{\hR}{h_{R}}
\newcommand{\Yd}{\Upsilon_{d}}
\newcommand{\Ys}{\Upsilon_{\ast}}
\newcommand{\Ydyn}{\Upsilon_{\rm dyn}}
\newcommand{\Tm}{T_{m}}
\newcommand{\so}{s_{0}}
\newcommand{\sigdyn}{\Sigma_{\rm dyn}}
\newcommand{\sigbar}{\Sigma_{\rm bar}}
\newcommand{\siglos}{\sigma_{\rm los}}
\newcommand{\rh}{r_{h}}
\newcommand{\LV}{L_{V}}
\newcommand{\epso}{\epsilon_{0}}
\newcommand{\kap}{\kappa}
\newcommand{\Mhalf}{M_{\rm half}}
\newcommand{\Mdyn}{M_{\rm dyn}}
\newcommand{\Mbar}{M_{\rm bar}}
\newcommand{\Mjmem}{M_{J,{\rm mem}}}
\newcommand{\Mjstd}{M_{J,{\rm std}}}
\newcommand{\fgas}{f_{\rm gas}}
\newcommand{\Sgal}{S_{\rm gal}}
\newcommand{\Gstr}{G_{\rm Strigari}}
\newcommand{\fp}{f_{p}}
\newcommand{\mustar}{\mu_{\ast}}
\newcommand{\AIC}{\Delta{\rm AIC}}
\newcommand{\chisq}{\chi^{2}}

% ---------------- title ------------------------------------------------------
\title{\bfseries Galactic Rotation Without Dark Matter\\
from Elastic Membrane Cosmology:\\[4pt]
\large The Geometric Mean Law and Observational Tests}
\author{Shinobu Sakaguchi\\[2pt]
\small Sakaguchi Seimensho (Independent Research), Shis\={o}, Hyogo, Japan\\
\small \texttt{sakaguchi-physics.com}}
\date{April 2026 (v4.7.8, 2026-04-21 affiliation rev.)}

% =============================================================================
\begin{document}
\maketitle
\thispagestyle{empty}

% ---------------- Abstract --------------------------------------------------
\begin{abstract}
\small\noindent
We present a membrane cosmology framework in which gravitational strength is modulated by the folding state
of a higher-dimensional elastic membrane, producing an effective gravitational acceleration
$\gobs = (\gN + \sqrt{\gN^2 + 4\gc\gN})/2$ that accounts for observed galactic rotation curves without invoking
dark matter. The central result is the \emph{geometric mean law}: the critical acceleration
$\gc = \eta\sqrt{\ao G \Sigma_0}$ is determined as the exact geometric mean ($\alpha=0.5$) of the universal
MOND scale $\ao$ and the local dynamical surface density $G\Sigma_0$, established from 175 SPARC galaxies with
$\AIC = -130$ versus MOND and confirmed by extended jackknife resampling. The Condition~15 final form
$\gc = 0.584\,\Yd^{-0.361}\sqrt{\ao\,\vflat^2/\hR}$ achieves $R^2 = 0.607$ with scatter $= 0.286$ dex,
established as an intrinsic limit (94\% irreducible). The membrane rigidity $\kap \sim 10^{-40}$ kpc$^2$ is
effectively zero, simplifying the Lagrangian to $\mathcal{L} \approx U(\epsilon; c)$ with no gradient term.
Each radius independently obeys $\mu(x) = x/(1+x)$ at RMS $= 0.090$ dex. SPARC-internal LITTLE THINGS dwarf
verification ($N=8$, $\vflat = 15\text{--}55$ km/s) yields scatter $= 0.232$ dex with KS $p=0.67$ versus the
full sample, confirming C15 consistency across the mass range. MOND rejection is established through three
independent methods ($p = 2 \times 10^{-27}$, $p = 8 \times 10^{-33}$, $p < 0.001$).

\smallskip\noindent
Version 4.7.8 extends the framework to pressure-supported dwarf spheroidal (dSph) galaxies. Direct C15
application to 31 dSph systems fails with $+1.5$ dex systematic bias, traced to an inversion of the
Jeans-mass hierarchy: $\Mjmem \ll \Mjstd$ in dSph (vs $\Mjstd \lesssim \Mjmem$ in SPARC), reversing the
direction of membrane--baryon coupling. In the $c \to 0$ limit of $U(\epsilon; c)$, the thermal Bernoulli
variance $\so(1-\so)$ predicts a universal acceleration scale $\Gstr = \so(1-\so)\,\ao = 0.228\,\ao$ with
$\Tm = \sqrt{6}$. The prediction is verified in two independent populations: 31 dSph galaxies ($0.240\,\ao$,
5\% agreement) and 30 radius points from four SPARC bridge galaxies with outer $\Ydyn > 10$ ($0.219\,\ao$,
4\% agreement). All four bridge galaxies (ESO444-G084, NGC2915, NGC1705, NGC3741) show statistically
significant radial increase of $\Ydyn$ ($z = 2.4$ to $21.6$), confirming a continuous
C15~$\to$~Strigari transition within individual galaxies. Multivariate regression
$\log\gobs$ vs $(\log\sigma, \log\rh, \log\LV)$ verifies $\LV$-independence at machine precision
($\beta_{L_V} = -2\e{-7}$, $R^2 = 1.0000$), confirmed independently by partial correlation. A four-component
scatter decomposition shows 89.4\% of the 0.350 dex residual scatter is physical, with the Milky Way and M31
satellite populations yielding identical mean offsets from the Bernoulli prediction ($+0.030$ vs $+0.032$
dex), indicating that $\Gstr$ is independent of the host galactic environment.

\smallskip\noindent\textbf{Keywords:} dark matter alternatives; galactic rotation curves; membrane cosmology;
geometric mean law; modified gravity; weak gravitational lensing; dwarf spheroidal galaxies; Strigari relation
\end{abstract}

\clearpage

% =============================================================================
\section{Introduction}
% =============================================================================

The flat rotation curves of spiral galaxies~\citep{RubinFord1980} and the success of the Tully--Fisher
relation~\citep{TullyFisher1977} have long been interpreted as evidence for dark matter halos. Alternative
approaches, including Modified Newtonian Dynamics (MOND;~\citealt{Milgrom1983}), Emergent
Gravity~\citep{Verlinde2017}, and various scalar-tensor theories, seek to reproduce these observations
through modifications of gravitational dynamics rather than new matter species.

Here we propose a membrane cosmology framework rooted in brane-world scenarios~\citep{RandallSundrum1999}.
The central idea is that gravitons propagate in a $(4+1)$-dimensional bulk, but are partially screened by
folds in the membrane carrying Standard Model matter. This screening modulates the effective Newton constant
locally, producing rotation-curve fits without a dark matter halo.

The key advance is the \emph{geometric mean law}: the critical acceleration $\gc$ that governs the transition
from Newtonian to MOND-like behaviour is not a universal constant (as MOND assumes) but varies from galaxy
to galaxy as the geometric mean of a universal scale $\ao$ and a local dynamical surface density $G\Sigma_0$.
This law is established from 175 SPARC galaxies with $\AIC = -130$ versus MOND. In v4.7, the theory is
substantially simplified: the membrane rigidity $\kap$ is confirmed to be zero, eliminating the gradient
term from the Lagrangian and reducing the free parameter count by four.

The paper is organised as follows. Section~2 presents the theoretical framework including the $\kap=0$
simplification. Section~3 establishes the geometric mean law and Condition~15 final form. Section~4
summarises statistical verification. Section~5 presents observational tests including weak lensing and dwarf
galaxy verification. Section~6 discusses comparison with MOND and remaining challenges. Section~7 extends
the framework to pressure-supported dwarf spheroidal systems via the $c \to 0$ limit of the elastic potential,
deriving the Bernoulli prediction $\Gstr = \so(1-\so)\,\ao$ and verifying continuous C15$\to$Strigari
transition in SPARC bridge galaxies, with $\LV$-independence verified at machine precision and scatter
decomposed into its physical components. Section~8 summarises.

% =============================================================================
\section{Theoretical Framework}
% =============================================================================

\subsection{Brane-world setup and gravitational screening}
Consider a $(4+1)$-dimensional spacetime with the Randall--Sundrum metric. Standard Model fields are confined
to the membrane at $y=0$. When the membrane is folded with local depth $D(r)$, gravitational flux lines are
partially redirected. We parametrise the fold by a dimensionless field $\epsilon(x)$ with equilibrium value
satisfying $U'(\epso) = 0$, where:
\begin{equation}
U(\epsilon; c) = -\epsilon - \frac{\epsilon^2}{2} - c\ln(1 - \epsilon), \qquad 0 < \epsilon < 1.
\label{eq:Uepsc}
\end{equation}
The effective gravitational acceleration felt by a test mass is:
\begin{equation}
\gobs(r) = \frac{\gN(r) + \sqrt{\gN(r)^2 + 4\gc\gN(r)}}{2}.
\label{eq:gobs}
\end{equation}
This is mathematically identical to the MOND simple interpolation function but arises from the equilibrium
condition of the elastic potential $U(\epsilon;c)$ --- a first-principles derivation rather than an empirical
ansatz. The general-$c$ equilibrium condition is:
\begin{equation}
\frac{c - 1 + \epsilon^2}{1 - \epsilon} = \frac{\gN}{\gc}.
\label{eq:eqcond}
\end{equation}

\subsection{Local equilibrium and membrane rigidity (v4.7)}
The membrane Lagrangian in principle contains a gradient term:
$\mathcal{L}[\epsilon] = U(\epsilon;c) + \kap (d\epsilon/dr)^2$. We estimate $\kap$ directly from SPARC
rotation curve data using both the Euler--Lagrange equation and the $E_{\rm grad}/E_{\rm local}$ ratio at each
radius.

\textbf{Result:} $\kap \sim 10^{-40}$ kpc$^2$ with 1.03 dex scatter and CV $= 3.14$ (non-universal).
The rigidity length $\ell_\kap \sim 0$ kpc: the membrane has no spatial elastic memory. The inverse
prediction $B = 2\kap/(U'' \hR^2) = 0.000$, compared to the empirical value $0.51$ from previous versions
--- a complete failure. Consequently:
\begin{enumerate}[leftmargin=*,itemsep=2pt]
\item The gradient term is physically irrelevant: $\mathcal{L} \approx U(\epsilon; c)$.
\item The membrane responds independently at each radius: each point on the rotation curve obeys
$\mu(x) = x/(1+x)$ with a single $\gc$ per galaxy. Simple interpolation applied independently at 3003 data
points across 167 galaxies gives RMS $= 0.090$ dex.
\item The Euler--Lagrange derivation pathway is abandoned.
\item Equations~\eqref{eq:Uepsc}--\eqref{eq:eqcond} constitute the complete dynamical content of the framework.
\end{enumerate}

% =============================================================================
\section{The Geometric Mean Law}
% =============================================================================

\subsection{Measurement of $\gc$}
For each SPARC galaxy, $\gc$ is measured directly from the Radial Acceleration Relation (RAR) without using
the scale radius $r_s$, avoiding circularity with the rotation curve fit. The characteristic acceleration
scale is $G\Sigma_0 = \vflat^2/\hR$.

\subsection{Core result}
\begin{equation}
\gc = \eta \cdot \ao^{(1-\alpha)} \cdot (G\Sigma_0)^{\alpha}, \qquad \alpha = 0.545 \pm 0.041.
\label{eq:geomean}
\end{equation}

\begin{center}\small
\begin{tabular}{@{}lrl@{}}
\toprule
Hypothesis & $p$-value & Decision \\
\midrule
$\alpha = 0$ (MOND: $\gc = $ const) & $2.0 \times 10^{-27}$ & Rejected \\
$\alpha = 0.5$ (geometric mean) & $0.27$ & Cannot reject \\
$\alpha = 1.0$ (pure $G\Sigma_0$) & $2.2 \times 10^{-21}$ & Rejected \\
\bottomrule
\end{tabular}
\end{center}

\subsection{Model comparison}
\begin{center}\small
\begin{tabular}{@{}lcr@{}}
\toprule
Model & Params & $\AIC$ vs MOND \\
\midrule
A: MOND ($\gc = \ao$) & 0 & 0 (baseline) \\
C: Geometric mean ($\alpha=0.5$ fixed) & 1 & $-130$ \\
B: Geometric mean ($\alpha$ free) & 2 & $-130$ \\
\bottomrule
\end{tabular}
\end{center}

\subsection{Condition~15 final form (v4.7 definitive)}
\begin{equation}
\gc = 0.584 \cdot \Yd^{-0.361} \cdot \sqrt{\ao \cdot \vflat^2/\hR}.
\label{eq:C15}
\end{equation}
\begin{center}\small
\begin{tabular}{@{}llp{60mm}@{}}
\toprule
Parameter & Value & Note \\
\midrule
$\alpha$ (fixed) & $0.5$ & Optimal; $\alpha$ free gives $+0.3\%$ improvement \\
$\beta$ ($\Yd$ exponent) & $-0.361 \pm 0.078$ & LOO verified, $\AIC = -14.2$ \\
$\eta_0$ & $0.584$ & (see Sec.~6.8) \\
$R^2$ & $0.607$ & Scatter $= 0.286$ dex (intrinsic limit) \\
LOO scatter & $0.290$ dex & No overfitting \\
$b_{\vflat}/(-b_{\hR})$ & $1.985$ & Theory prediction: 2.0 (exact match) \\
MOND rejection & $p = 1.66 \times 10^{-53}$ & Decisive \\
\bottomrule
\end{tabular}
\end{center}

\textbf{Scatter decomposition:}
\begin{center}\small
\begin{tabular}{@{}lcr@{}}
\toprule
Component & $\sigma$ (dex) & Fraction \\
\midrule
Intrinsic physical variance & $0.227$ & $63.5\%$ \\
$\gc$ measurement error & $0.152$ & $28.3\%$ \\
$\Yd + \vflat + \hR$ input & $\sim 0.08$ & $8.2\%$ \\
\textbf{Total} & $0.286$ & $100\%$ \\
\bottomrule
\end{tabular}
\end{center}

\subsection{Jackknife verification}
\begin{center}\small
\begin{tabular}{@{}lp{80mm}@{}}
\toprule
Test & Result \\
\midrule
Half-split $\times 1000$ & $\alpha_{\rm train} = 0.545 \pm 0.037$, $\alpha_{\rm test} = 0.550 \pm 0.038$ \\
Leave-50-out $\times 500$ & $\alpha = 0.546 \pm 0.024$, range $[0.483, 0.638]$ \\
Bootstrap $\times 1000$ & $\alpha = 0.548 \pm 0.038$, 95\% CI $[0.478, 0.626]$ \\
Improves over MOND & $100\%$ of splits (median $31\%$) \\
\bottomrule
\end{tabular}
\end{center}

\subsection{Decomposition and origin of $\alpha = 0.5$}
Element substitution analysis reveals that the $\gc$ measurement method (tanh fit vs direct median) is
irrelevant; what matters is per-galaxy $\Yd$ optimisation. Direct median with $\Yd = 0.5$ fixed gives
$\alpha = 0.432$, but with $\Yd$ individually optimised gives $\alpha = 0.546$ --- identical to the tanh result.
The origin of $\alpha = 0.5$ is identified as an algebraic consequence of the deep-MOND limit. With
$\gN \ll \gc$, Eq.~\eqref{eq:gobs} reduces to $g_{\rm eff} = \sqrt{\gN\gc}$. Direct fitting yields
$\alpha = 0.471 \pm 0.036$, $p(0.5) = 0.43$ (A-grade). Three-layer decomposition: deep-MOND (0.47) +
MOND transition ($+0.05$) + residual scatter ($+0.02$) $= 0.545$.

\subsection{Component decomposition}
$\gc \sim \vflat^{1.10}\,\hR^{-0.56}$, $R^2 = 0.529$. Both exponents are consistent with the geometric mean
prediction ($\vflat^{1.0}\,\hR^{-0.5}$): $p(\vflat = 1.0) = 0.24$, $p(\hR = -0.5) = 0.49$. Total baryonic mass
$\Mbar$ alone gives $R^2 = 0.021$ --- mass does not drive $\gc$. The surface density proxy
$\sigdyn = \vflat^2/\hR$ is the essential quantity.

\subsection{C15 as minimal sufficient statistic (v4.7)}
Six models incorporating a galaxy-level plasticity indicator $\Sgal$ (constructed from $\fgas$, $\Yd$,
SBdisk0, T-type) are tested against C15. All six are rejected (LOO improvement $< 2\%$, no model satisfies
both LOO $> 5\%$ and $\AIC < -5$). The decisive finding: in model M1 (C15 + $\Sgal$), the $\Sgal$
coefficient reverses sign ($+0.44$) compared to the raw partial correlation ($\rho = -0.23$), indicating
that $\Sgal$'s apparent predictive power is entirely absorbed by $\sigdyn$ already present in C15. Forward
stepwise regression with all available parameters ($I_0$, $\fgas$, T-type, quality $Q$) achieves at most
$1.7\%$ improvement (LOO). Distance, inclination, and velocity errors are non-significant ($p > 0.3$). C15
is the minimal sufficient statistic for $\gc$ prediction. Grade: A.

% =============================================================================
\section{Statistical Verification}
% =============================================================================

\subsection{SPARC 167 galaxies}
Extending $\gc$ measurement to 167 SPARC galaxies (quality cut $\chisq/\text{dof} < 10$):
median $\gc/\ao = 0.24$, with systematic Hubble-type dependence (Spearman $r = -0.760$, $p < 0.0001$).
$\gc$ decreases by a factor of 22 from Sab to Im type. The null correlation of $r_s$ with $\gc$
($r = 0.018$, n.s.) demonstrates that $\gc$ is determined by membrane material properties, not spatial
location.

\subsection{First-principles derivation of $F_{\rm conf}$}
The confining term $F_{\rm conf} = -c\ln(1 - \epsilon)$ in $U(\epsilon; c)$ is derived from the configuration
entropy of the kink--fold network via forest~$+$~cavity mean-field effective theory (Lemmas 1--4, all
established). Observational confirmation: $z_{\rm kin}$ shows $r = -0.507$ ($p = 2 \e{-12}$) anti-correlation
with $\gc/\ao$, and two independent proxies agree at $r = 0.809$ ($p < 10^{-38}$), supporting the physical
reality of the fold density $n_{\rm fold}$. Self-consistency: $\alpha = 0.5$ follows from the deep-MOND
algebraic structure and the BTFR--$\vflat$--$\hR$ closure (slope $= 0.545 \pm 0.075$, prediction 0.5,
$0.6\sigma$).

% =============================================================================
\section{Observational Tests}
% =============================================================================

\subsection{Weak-lensing Q-value test}
Nine galaxy clusters from the~\citet{Miyaoka2019} sample with HSC-SSP weak-lensing data: $Q > 1$ in 9/9
clusters (100\%). Wilcoxon $p = 0.0020$ (Grade A). Sign of $\alpha$ confirmed at 98.4\%.

\subsection{HSC-SSP cl1 analysis}
\begin{center}\small
\begin{tabular}{@{}lcc@{}}
\toprule
Model & $\chisq/\text{dof}$ & $\AIC$ vs NFW \\
\midrule
NFW ($M_{200}, c_{200}$) & $1.98$ & $0$ \\
MOND ($\gc = \ao$, Abel) & $1.69$ & $-4.9$ \\
Membrane ($\gc$ free, Abel) & $1.54$ & $-4.4$ \\
\bottomrule
\end{tabular}
\end{center}
$\gc = 1.58\,\ao$ (68\% CI: $[1.12, 2.00]$). $\ao$ is within 95\% CI. Level B$+$ (systematic issues:
$M_{\rm lens}/M_{\sigma_v} = 7.5$, $\gc$ method dependence). Future surveys (Rubin LSST, Euclid) needed.

\subsection{LITTLE THINGS dwarf verification (v4.7)}
Eight LITTLE THINGS dwarf irregular galaxies~\citep{Oh2015} are identified within the SPARC 175 sample,
spanning $\vflat = 15\text{--}55$ km/s and $\gc/\ao = 0.1\text{--}0.65$:

\begin{center}\small
\begin{tabular}{@{}lccc@{}}
\toprule
Metric & LT subset (N=8) & Full SPARC & non-LT \\
\midrule
C15 scatter (dex) & $0.232$ & $0.286$ & $0.288$ \\
Median bias (dex) & $-0.189$ & $+0.000$ & $+0.004$ \\
KS test $p$ & $0.67$ & --- & --- \\
\bottomrule
\end{tabular}
\end{center}

KS $p = 0.67$ confirms that dwarf residuals are drawn from the same distribution. C15 holds uniformly across
the mass range. Combined with the full-sample LITTLE THINGS result ($N = 178$, $\alpha = 0.576 \pm 0.047$,
$p(0.5) = 0.109$), three independent MOND rejection pathways are established.

\paragraph{Non-SPARC external verification (v4.7.7).}
To extend the dwarf verification beyond the SPARC sample, we applied C15 to 18 non-SPARC LITTLE THINGS
galaxies from~\citet{Oh2015} with baryon-separated rotation curves from VLA HI and Spitzer 3.6~$\mu$m
photometry. Galaxy-specific $\Yd$ values were estimated from SED-derived stellar masses~\citep{Zhang2012}
and V-band luminosities. After quality cuts excluding 5 galaxies with known kinematic anomalies (starburst,
outflow contamination, or insufficient $V_{\rm bar}$ coverage), the remaining 11 high-quality galaxies yield
scatter $= 0.289$ dex with bias $= -0.004$ dex. A KS test confirms that non-SPARC and SPARC-internal
residuals are drawn from the same distribution ($p = 0.73$). Combined with the SPARC-internal result ($N =
8$, scatter $= 0.232$ dex), C15 is verified across the full dwarf mass range $\vflat = 12\text{--}125$ km/s
with no systematic deviation. Grade: A (upgraded from B$+$ via independent external confirmation).

% =============================================================================
\section{Comparison with MOND and Remaining Challenges}
% =============================================================================

\subsection{MOND as special case}
MOND's $\ao$ is recovered as the all-galaxy average of $\gc$. An ``average galaxy'' ($G\Sigma_0 \sim \ao/\eta^2$)
gives $\gc \sim \ao$, recovering MOND. MOND is a special case of the membrane framework. MOND rejection is
established through three independent methods: RAR-based ($p = 2\e{-27}$), deep-MOND inversion
($p = 8\e{-33}$), and LITTLE THINGS independent ($p < 0.001$).

\subsection{$V_{\rm bar}$ bottleneck}
The primary obstacle to Level A external verification is the $V_{\rm bar}$ bottleneck: baryon-separated
rotation curves are publicly available only in SPARC and LITTLE THINGS. PROBES (3{,}158 galaxies) and MaNGA
DynPop (10{,}296) provide only $V_{\rm obs}$, yielding $\alpha \sim 1.0$ with fixed $\Yd = 0.5$. These are
repositioned for large-sample precision measurement rather than independent verification. BIG-SPARC ($\sim
4000$ galaxies, in preparation) will provide the definitive external test.

\subsection{$\kap = 0$ and theory simplification (v4.7)}
The $\kap = 0$ result (Sec.~2.2) has far-reaching consequences for the theory structure:
\begin{center}\small
\begin{tabular}{@{}lll@{}}
\toprule
Aspect & Before (v4.6) & After (v4.7) \\
\midrule
Lagrangian & $\mathcal{L} = U(\epsilon;c) + \kap(d\epsilon/dr)^2$ & $\mathcal{L} \approx U(\epsilon;c)$ \\
$\gc$ structure & C15 global + C14 radial & C15 only \\
Scatter & 0.286 dex (improvable?) & 0.286 dex (intrinsic limit) \\
Free params removed & A, B, $\tau$, $\kap$ & All four eliminated \\
\bottomrule
\end{tabular}
\end{center}

The Condition~14 dynamical formulation $g_{c,\rm eff}(r) = \gc[0.12 + 0.51\exp(-r/(4.7\hR))]$ reported in
v4.6 is retracted. The pattern is a projection of C15 prediction bias onto the radial coordinate. The
conceptual hypothesis of plastic/elastic two-phase membrane structure (M3 $\epsilon$ threshold,
$\AIC_c = -5.1$) remains at B$-$grade as an interpretive framework.

\subsection{$\hR$ partial correlation}
The scale-length $\hR$ shows zero raw correlation with $\gc$ ($\rho = -0.022$) but a significant partial
correlation after controlling for $\vflat$ ($\rho = -0.472$). Decomposition: 69\% from outer-disk MOND
transition (A-grade), 18\% from $x(r)$ profile higher-order effects (B$-$grade, reattributed from gradient
energy in v4.7), 13\% residual.

\subsection{BTFR slope and Simpson's paradox}
The membrane BTFR slope $= 1.20$ is a statistical confound independent of $\gc$: setting $\gc = \ao$
(constant) yields the same slope 1.196. Simpson's paradox arises from $\vflat$ driving both BTFR axes.
Within-bin slope $= 0.542 \pm 0.069$ (bootstrap 95\% CI $[0.400, 0.668]$, $p(0.5) = 0.494$) provides a
non-circular confirmation of $\alpha = 0.5$. In dwarf galaxies ($\vflat < 47$ km/s), $b_{gc} = 2.32$: the
membrane fold state dominates rotation velocity, a distinctive prediction with no analogue in standard dark
matter models.

\subsection{Deep-MOND $\gc$--$\sigbar$ independence}
In the deep-MOND regime (bottom 30\% by $\gc$), $\gc$ is uncorrelated with baryonic surface density:
$r(\gc, \sigbar) = +0.066$. In contrast, $r(\gc, \sigdyn) = +0.727$. The ratio $\sigbar/\sigdyn$
anti-correlates with $\gc$ ($r = -0.648$): from $\sim 0.7$ in low-$\gc$ galaxies to $\sim 0.14$ in high-$\gc$.
The membrane responds to dynamical stress-energy, not baryonic mass content. The self-referential closure
$\gc \sim \sigbar^{1/3}$ fails because its premise is empirically false. Grade: A.

\subsection{Intrinsic scatter limit}
Two independent optimisation pathways fail to reduce the 0.286 dex scatter: (i) $\alpha(\gc)$ tanh transition
achieves 3\% improvement ($\AIC = -2.4$, 4 extra parameters --- insufficient); (ii) all additional galaxy
parameters achieve at most 1.7\% (LOO). The scatter-minimising $\alpha(\gc)$ parameters contradict the
sliding-window $\alpha$ in direction, confirming that local statistical slopes and global prediction accuracy
are distinct quantities. The Rotmod-derived local $\gc$ (0.244 dex) provides a lower bound. Grade: A.

\subsection{Plasticity narrative and self-consistent equation (v4.7.6)}
The $\Yd$ dependence of $\eta$ is a structural consequence of the BTFR + GML system, not a ``plasticity
proxy''. All plasticity indicators ($\fgas$, SBdisk0, T-type) are different aspects of $\sigdyn$ variation
already captured by C15. The morphological-type bridge (T-type vs $\Sgal$, $\rho = +0.89$,
$p = 4\e{-60}$) establishes a strong conceptual correspondence but adds no predictive power.
``Plasticity'' is an explanatory narrative, not a predictive variable. Grade: A (bridge) / X ($\Sgal$
prediction).

\paragraph{Theoretical basis for $\eta_0$ (v4.7.6).}
The C15 normalisation $\eta_0 = 0.584$ can be given a partial first-principles interpretation through the
membrane self-consistent equation. Define the elastic fraction $s = 1 - \fp$, where $\fp$ is the plastic
content of the membrane at a given galaxy. The dimensionless dynamical structure parameter
$Q = \sqrt{\vflat^2/(\ao \hR)}$ measures the fold-driving stress. The equilibrium elastic fraction satisfies:
\begin{equation}
s = \frac{1}{1 + \exp(-\Delta U(sQ)/\Tm)},
\label{eq:SCE}
\end{equation}
where $\Delta U(c) = -3/2 + 2\sqrt{c} - c/2 - (c/2)\ln(c)$ for $0 < c \le 1$ is the potential energy
difference between the folded and unfolded membrane states, derived from $U(\epsilon; c)$ (Eq.~\ref{eq:Uepsc}).
The effective membrane temperature $\Tm = \sqrt{6}$ is the $Z_2$ symmetry-breaking critical temperature
(Lemma~5). Key properties of Eq.~\eqref{eq:SCE}: $\Delta U(0) = -3/2$ (deep plastic well), $\Delta U(1) = 0$
with $\Delta U'(1) = 0$ (inflection point), guaranteeing a unique solution for each $Q$.

Applying Eq.~\eqref{eq:SCE} to the 159 SPARC galaxies in the TA3 sample, the SPARC-averaged elastic fraction
is $\langle s \rangle = 0.58$, matching $\eta_0 = 0.584$ to within 1\%. This identifies $\eta_0$ as the mean
elastic fraction of the membrane across the observed galaxy population: $\eta_0 = \langle 1 - \fp
\rangle_{\rm SPARC}$. Grade: B$+$.

\paragraph{$\Yd$ as independent variable (v4.7.6).}
The extended model $\gc/\ao = A\,s(Q, \Tm)\,Q\,\Yd^{\beta}$ with $\Tm = \sqrt{6}$ yields $\beta = -0.331$ and
scatter $= 0.2845$ dex, compared to the C15 empirical $\beta = -0.361$ (8\% difference). Crucially, the
residuals of the pure $s(Q)$ model ($\Yd^{\beta}$ omitted) correlate with $\Yd$ at $r = -0.296$
($p = 7 \e{-5}$), demonstrating that $\Yd$ cannot be derived from $Q$ alone. The membrane state space is at
least two-dimensional: dynamical structure ($Q$) and material composition ($\Yd$). Grade: A ($\Yd$
independence).

\paragraph{Partial correlation test (v4.7.7).}
To test whether the $\gc$--$\Yd$ correlation is mediated by gas fraction, we compute the partial correlation
$r(\log\gc, \log\Yd | \fgas)$ using the SPARC 175 sample ($\fgas = 1.33\,M_{\rm HI}/(1.33\,M_{\rm HI} + \Yd
L_{3.6})$). The partial correlation is $r = -0.437$ ($p = 1.6\e{-9}$), slightly stronger than the raw
$r = -0.428$. The attenuation ratio is 1.02: conditioning on $\fgas$ does not weaken the correlation. The
decisive result is $r(\gc, \fgas) = +0.061$ ($p = 0.42$): $\gc$ is uncorrelated with gas fraction. This rules
out the common-cause model $\fp \to \gc$ and $\fp \to \Yd$ mediated by $\fgas$. Instead, $\Yd$ itself is the
direct observational proxy of the membrane plastic fraction $\fp$. Grade: A.

\paragraph{Resolution of the $\beta$ discrepancy (v4.7.7).}
The 8.3\% difference between the self-consistent equation prediction ($\beta = -0.331$ at $\Tm = \sqrt{6}$)
and the C15 empirical value ($\beta = -0.361$) is fully explained as a finite-temperature absorption effect.
A systematic $\Tm$ scan confirms that at $\Tm \to \infty$ ($s = 0.5 =$ const), the model recovers $\beta =
-0.361$ exactly: C15 implicitly assumes constant $s$. At finite $\Tm = \sqrt{6}$, the variation of $s(Q)$
across the galaxy population absorbs 8.3\% of the $\Yd$ dependence, reducing the direct $\beta$ from $-0.361$
to $-0.331$. This absorption requires no parameter adjustment and arises from the $Q$-dependent structure of
$\Delta U(c)$. The scatter-minimising $\Tm = 2.35$ is consistent with $\sqrt{6} = 2.449$ to within
$0.24\sigma$, providing independent empirical support for the $Z_2$ SSB critical temperature. Grade: A
($\beta$ discrepancy resolved; $\Tm = \sqrt{6}$ confirmed by scatter optimisation).

\subsection{Weak Lensing RAR: KiDS-1000 and HSC-SSP Y3 (v4.7.6)}
We examined the KiDS-1000 weak lensing Radial Acceleration Relation~\citep{Brouwer2021} for consistency with
C15. The KiDS lensing RAR extends the rotation-curve RAR by two decades in $\gobs$ into the low-acceleration
regime ($\gobs \sim 10^{-15}$ to $10^{-12}$ m/s$^2$), using stacked Excess Surface Density (ESD) profiles of
$\sim 1$ million isolated galaxies from the KiDS-bright sample.

Following~\citet{Brouwer2021}, we adopt their nominal circumgalactic gas (CGM) model
($M_{\rm hot}/M_\ast = 1$, isothermal profile, truncation radius $R_{\rm acc} = 100$ kpc). A universal
correction factor $C(\gbar) = {\rm ESD}_{\rm hotgas}(\gbar)/{\rm ESD}_{\rm no\text{-}hotgas}(\gbar)$ is derived
from the full-sample (no mass bins) ESD pair and applied to four stellar mass bins
($\log_{10} M_\ast/M_\odot = [8.5, 10.3, 10.6, 10.8, 11.0]$). Per-bin $\gc$ is then fitted using the membrane RAR.

Under the nominal CGM model, the hot-gas-corrected per-bin $\gc$ values show a positive $M_\ast$ slope
($+0.138$ in $\log\gc$ vs $\log M_\ast$), consistent in sign and order of magnitude with the C15 prediction
($+0.075$). The per-bin ratio $\gc^{\rm obs}/\gc^{\rm C15}$ ranges from 0.77 to 1.05 across all four bins,
with Bin~4 ($\log M_\ast = 10.9$) achieving exact agreement (ratio $= 1.00$).

However, a systematic sensitivity analysis reveals that this result is not robust to CGM mass uncertainty.
Scanning $M_{\rm hot}/M_\ast$ over the observationally permitted range $[0.3, 3.0]$~\citep{Tumlinson2017,
Babyk2018}, the formal $\Delta\chisq({\rm MOND} - {\rm C15})$ computed with the full $60 \times 60$
covariance matrix reverses sign: $\Delta\chisq = -10.6$ at $\fgas = 0.3$ (MOND preferred) versus $+147.7$ at
$\fgas = 1.0$ (C15 preferred). Introducing a physically motivated $M_\ast$-dependent correction
($R_{\rm acc} \propto M_\ast^{0.4}$) further reduces the C15 preference to $\Delta\chisq = +0.8$ at
$\fgas = 1.0$, rendering the two models statistically indistinguishable.

\begin{sloppypar}
The~\citet{Brouwer2021} colour and S\'ersic-index bins (their Fig.~8) provide a complementary test.
Splitting by colour or S\'ersic index, we find $\gc({\rm red})/\gc({\rm blue}) = 2.3$ to $2.8$ and
$\gc({\rm high-}n)/\gc({\rm low-}n) = 2.4$ to $2.5$, with the ordering stable across $\fgas = 0$ to 1.0
($\Delta\chisq({\rm uniform} - {\rm split}) = +127$ for colour, $+67$ for S\'ersic). This result is
CGM-robust because the hot-gas correction affects both morphological bins approximately equally, preserving
the $\gc$ ratio.
\end{sloppypar}

However, a SPARC cross-check reveals that this ordering is dominated by the colour-magnitude relation rather
than a pure morphological effect. In SPARC, the median $\gc$ ratio between early-type ($T \le 3$) and
late-type ($T > 3$) galaxies is 0.92, with no statistically significant difference (Mann--Whitney $p = 0.35$).
C15 predicts a ratio of 1.7, arising entirely from $\sigdyn = \vflat^2/\hR$ differences, with zero
contribution from $\Yd$ (both groups share median $\Yd = 0.30$ in the TA3 sample). The KiDS ratio of 2.3 to
2.8 exceeds both the SPARC observation and the C15 prediction, indicating that the colour/S\'ersic bins
primarily separate galaxies by stellar mass rather than by morphological type, with the $\gc$ ordering
reflecting the underlying $M_\ast$--$\sigdyn$ correlation.

In summary, the KiDS lensing data are consistent with both C15 and MOND at the current level of systematic
uncertainty. The stellar-mass test cannot discriminate between the two models due to CGM mass uncertainty,
and the colour/S\'ersic test, while CGM-robust, is dominated by the colour-magnitude relation rather than an
independent morphological signal. Definitive tests will require future surveys combining high-precision weak
lensing (Euclid, Rubin LSST) with resolved CGM measurements (eROSITA, SKA), and colour-binned analysis
within narrow stellar-mass slices to isolate pure morphological effects. Grade: B (both tests). $M_\ast$
test: CGM-limited. Colour/S\'ersic test: CGM-robust but colour-$M_\ast$ degenerate.

\paragraph{HSC-SSP Y3 independent lensing RAR (v4.7.4).}
To provide an independent test of the lensing RAR obtained with KiDS-1000, we constructed a fully
independent weak-lensing pipeline using the Hyper Suprime-Cam Subaru Strategic Program Year 3 (HSC-SSP Y3)
shape catalogue as source galaxies and the Galaxy and Mass Assembly Data Release 4 (GAMA DR4) spectroscopic
survey as lens galaxies. The HSC-SSP Y3 catalogue contains 35.7 million galaxies with shapes measured via
the Re-Gaussianization (regauss) method, reaching an effective source number density of 19.9 arcmin$^{-2}$
with a median $i$-band seeing of 0.6~arcsec. GAMA~DR4 provides accurate spectroscopic redshifts
($N_Q \ge 3$) and stellar masses across three equatorial fields (G09, G12, G15), each spanning 60~deg$^2$.

We selected isolated lens galaxies from GAMA~DR4 via the G3C group catalogue, retaining only central
galaxies (RankIterCen $= 1$) with $0.05 < z < 0.5$ and $\log_{10}(M_\ast/M_\odot) < 11.0$, the latter
matching the isolation criterion of~\citet{Brouwer2021}. This yielded 157{,}338 isolated lenses across three
fields: 49{,}272 (G09), 55{,}402 (G12), 52{,}664 (G15). The baryonic acceleration was computed per lens as
$\gbar = G M_{\rm gal}/R^2$, treating each lens as a baryonic point mass~\citep[][Eq.~2]{Brouwer2021}, with
pairs binned in 15 logarithmic $\gbar$ bins spanning $10^{-15}$ to $5\e{-12}$ m~s$^{-2}$. A total of 503
million lens--source pairs were accumulated.

Systematic corrections were applied as follows. The HSC regauss shear responsivity
$\mathcal{R} = 1 - \langle e_{\rm rms}^2 \rangle$ was computed per field ($\mathcal{R} \sim 0.86$). The
multiplicative shear bias for HSC Y3 regauss is controlled at $|m| < 10^{-2}$ through image-simulation
calibration~\citep{Li2022, Mandelbaum2018}, contributing $< 1\%$ ESD uncertainty. The observed acceleration
was obtained via the SIS approximation $\gobs = 4G\Delta\Sigma$~\citep[][Eq.~7]{Brouwer2021}. The cross
component ${\rm ESD}_\times$ is consistent with zero across all bins, confirming the absence of significant
additive systematics. Statistical uncertainties were estimated using a 24-patch spatial jackknife (8 per
field) yielding the full $15 \times 15$ covariance matrix.

A critical systematic-robustness test is the consistency of results across independent survey fields.
Fitting the interpolation $\gobs = \gbar/(1 - \exp(-\sqrt{\gbar/\gc}))$ independently in each field yields
$\gc = 2.91 \pm 0.19\,\ao$ (G09), $2.64 \pm 0.16\,\ao$ (G12), $2.64 \pm 0.23\,\ao$ (G15). The inter-field
consistency $\chisq = 1.37$ (2 d.o.f., $p = 0.50$) confirms that all three fields are statistically
consistent at the $1\sigma$ level, with no evidence for field-dependent systematics.

Combining all three fields yields $\gc = 2.73 \pm 0.11\,\ao$ for the full isolated lens sample. Comparing
C15 ($\gc$ free) to MOND ($\gc = \ao$ fixed) using the full jackknife covariance matrix yields
$\chisq_{\rm MOND} = 539.3$ versus $\chisq_{\rm C15} = 65.0$, giving $\Delta\chisq = +474$ and
$\AIC = +472$, corresponding to a $\sim 22$-sigma preference for C15 over MOND. The lensing RAR of isolated
galaxies at Mpc scales is incompatible with a universal acceleration scale $\gc = \ao$ as predicted by MOND.

The $\gc$--$M_\ast$ relation provides a further discriminant. Fitting $\gc$ in four stellar-mass bins
between $\log_{10} M_\ast = 8.5$ and $11.0$, we find a monotonic increase from $\gc = 2.00 \pm 0.27\,\ao$
(lowest $M_\ast$) to $3.41 \pm 0.21\,\ao$ (highest $M_\ast$), with a log--log slope of $+0.166 \pm 0.041$.
MOND predicts zero slope and is rejected at $4.0\sigma$ ($p < 0.001$). C15, which predicts a positive slope
of $+0.075$ from the $\sigdyn$ dependence, is consistent at the $2.2\sigma$ level ($p = 0.029$). The observed
slope being approximately twice the C15 prediction suggests that the $\sigdyn$ dependence at Mpc lensing
scales may be steeper than at kpc rotation-curve scales, potentially reflecting the transition from
baryon-dominated to halo-dominated regimes. This is a quantitative prediction for future refinement of the
membrane framework.

In summary, the HSC-SSP Y3 lensing RAR independently confirms the C15 framework at high significance:
(i) $\gc = 2.73 \pm 0.11\,\ao$ is inconsistent with the MOND universal value $\ao$ at $> 15\sigma$;
(ii) the $\gc$--$M_\ast$ slope rejects MOND at $4\sigma$ while remaining consistent with C15's predicted
positive slope;
(iii) all three GAMA fields yield consistent $\gc$ values ($p = 0.50$);
(iv) combined with the KiDS-1000 analysis, the lensing RAR evidence upgrades from B grade (KiDS alone,
CGM-limited) to A grade (independent HSC confirmation with systematic verification). The observed factor
$\sim 2$ excess of the slope over the C15 prediction is noted as a quantitative tension requiring
theoretical investigation, but does not affect the qualitative conclusion that $\gc$ varies with galaxy
properties as predicted by membrane cosmology and contrary to MOND.

\paragraph{Slope budget analysis (v4.7.6).}
The factor $\sim 2$ excess of the observed $\gc$--$M_\ast$ slope ($+0.166 \pm 0.041$) over the C15
prediction ($+0.075$) reported above constituted a $2.2\sigma$ tension. We decompose this excess into
identified physical contributions.

First, the 2-halo contribution to the lensing ESD slope was measured empirically by separating the 1-halo
and total (1h+2h) regimes in each $M_\ast$ bin. Fitting $\gc$ using only the 1-halo-dominated radial range
($R < 200$ kpc) and comparing to the full-range fit yields $\Delta\text{slope}_{\rm 2h} = +0.034 \pm 0.064$
($0.5\sigma$). The sign is consistent with the physical expectation that more massive galaxies reside in
denser large-scale environments, enhancing the outer ESD and steepening the apparent $\gc$--$M_\ast$
relation.

Second, the light-propagation model (Eqs.~10a--10d) predicts a small additional slope contribution of
$+0.017$ from the $R$-dependent gravitational Aharonov--Bohm phase shift (Sec.~2).

The combined slope budget is:
\begin{center}\small
\begin{tabular}{@{}lll@{}}
\toprule
Source & Slope contribution & Status \\
\midrule
C15 ($\sigdyn$ dependence) & $+0.075$ & A-grade (SPARC) \\
Model B (AB phase shift) & $+0.017$ & B-grade \\
2-halo environment & $+0.034 \pm 0.064$ & B-grade ($0.5\sigma$) \\
\midrule
Total predicted & $+0.126$ & \\
Observed & $+0.166 \pm 0.041$ & \\
Residual tension & $0.98\sigma$ & Resolved \\
\bottomrule
\end{tabular}
\end{center}

The slope tension is reduced from $2.2\sigma$ to $0.98\sigma$, within the expected statistical fluctuation.
While the individual 2-halo and Model~B contributions are each non-significant, their combined effect
accounts for the observed excess. Future surveys with higher signal-to-noise (Euclid Wide, Rubin LSST) are
expected to detect the 2-halo contribution at $2\text{--}3\sigma$ significance.

% =============================================================================
\section{Pressure-Supported Dwarf Spheroidal Extension}
% =============================================================================

The v4.7.7 framework was developed against rotation-supported systems (SPARC spirals and dwarf irregulars).
This section extends the framework to pressure-supported dwarf spheroidal (dSph) galaxies, exposing a
qualitatively new regime in which the baryonic driving of the membrane fails and the spontaneous $c \to 0$
equilibrium of $U(\epsilon; c)$ dominates. The resulting Bernoulli-variance prediction
$\Gstr = \so(1 - \so)\,\ao$ is verified in two statistically independent populations to 4--5\% accuracy.

\subsection{dSph sample construction}
We assembled a catalogue of 31 dSph galaxies: 15 Milky Way satellites ($\siglos$ from~\citealt{Walker2009};
\citealt{SimonGeha2007}; \citealt{Simon2019}), 11 M31 satellites~\citep{Collins2013, Tollerud2012}, 3
isolated systems (Cetus, Phoenix, Tucana), and 2 unclassified dwarfs. Half-light radii $\rh$ were taken
from~\citet{McConnachie2012}, with the Walker estimator $\Mdyn = 2.5\,\siglos^2\,\rh/G$ used for $\sigma$
back-calculation in the VizieR catalogue (coefficient 2.5, confirmed by
$\sigma_{\rm wolf}/\sigma_{\rm lit} = 0.997$). Dynamical accelerations were computed via the~\citet{Wolf2010}
relation $\gobs = 3\siglos^2/\rh$, and baryonic accelerations from stellar mass with $\Ys = 1$ assumed.

\subsection{C15 direct application fails: J3 regime inversion}
Applying the C15 formula (Eq.~\ref{eq:C15}) to the dSph sample yields a systematic offset of $+1.5$ dex with
scatter 0.9 dex. Four model variants (C15\_naive, SCE\_naive, SCE\_Q\_J, C15\_Q\_J) give statistically
indistinguishable residuals. The $\Yd$ transformation does not absorb the offset (scatter invariant under
$\Yd$ scan). Auxiliary hypotheses contribute insufficiently: H1 (3D topology) $+0.15$ dex, H2 (virial
coefficient $\beta_c = 2$ vs $\sqrt{3}$) $+0.25$ dex, combined $+0.40$ dex (29\% of the observed offset).
H3 (membrane hysteresis alone) is rejected by three independent tests (isolated $>$ satellite $\gc$; zero
distance correlation; inferior AIC versus constant-model).

The failure traces to an inversion of the Jeans-mass hierarchy. Define the membrane Jeans mass
$\Mjmem = (4\pi/3)\rho(c_s^2/g_{\rm eff})^3$ and the standard Jeans mass $\Mjstd$. In SPARC, $\Mjstd \le
\Mjmem$ and baryonic distribution drives the membrane state (C15 valid). In dSph, $\Mjmem \ll \Mjstd$ in 28
of 31 systems, and the membrane sets the structure while baryons follow. The direction of causality is
reversed, and $\gc$ is no longer a function of baryonic surface density $\Sigma_0$. Grade: A.

\subsection{Unified RAR: SPARC $\times$ dSph and J3 transition band}
Overlaying 3{,}389 SPARC RAR-cloud points from Rotmod\_LTG~\citep{Lelli2016} with the 31 dSph sample on the
$\log(\gbar/\ao)$--$\log(\gobs/\ao)$ plane reveals two populations occupying distinct
$\Ydyn = \gobs/\gbar$ ranges:

\begin{center}\small
\begin{tabular}{@{}lrrl@{}}
\toprule
Sample & $N$ & median $\Ydyn$ & Interpretation \\
\midrule
SPARC RAR cloud (all radii) & 3{,}389 & 2.75 & baryon-dominated \\
SPARC galaxy-level (outer) & 175 & 3 (typical) & C15 regime \\
Bridge galaxy outer (40\%) & 30 & 12.76 & J3 transition band \\
dSph & 31 & 35.11 & membrane-dominated \\
\bottomrule
\end{tabular}
\end{center}

We designate the band $10 < \Ydyn < 30$ the J3 transition band. Only 4 of 175 SPARC galaxies reach outer
$\Ydyn > 10$: ESO444-G084, NGC2915, NGC1705, NGC3741. These ``bridge galaxies'' span the SPARC--dSph gap and
permit radius-resolved verification of the transition.

\subsection{Reformulation via $c \to 0$ limit of $U(\epsilon; c)$}
In the J3 regime the coupling $c$ (set by baryonic driving) approaches zero, and the membrane enters its
intrinsic thermal equilibrium. The self-consistent equation Eq.~\eqref{eq:SCE} in the $Q \to 0$ limit yields
a universal spontaneous elastic fraction:
\begin{equation}
\so = \frac{1}{1 + \exp(3/(2\Tm))}.
\label{eq:s0}
\end{equation}
With $\Tm = \sqrt{6}$, this gives $\so = 0.3515$. The leading small-$Q$ correction is non-analytic:
\begin{equation}
s(Q) \approx \so + \alpha_s \sqrt{Q} + O(Q), \qquad \alpha_s = \frac{2E}{F\Tm}\so^{3/2} = 0.110,
\label{eq:sQ}
\end{equation}
where $E = \exp(3/(2\Tm))$ and $F = 1 + E$. The $\sqrt{Q}$ structure originates from the $\ln(1 - \epsilon)$
term in $U(\epsilon; c)$ and is characteristic of the $c \to 0$ singularity. Numerical SCE verification
confirms agreement with the asymptotic expansion to within 4\% relative error for $Q < 0.3$.

\subsection{Bernoulli prediction: $\Gstr = \so(1 - \so)\ao$}
In the $c \to 0$ regime the membrane occupies a thermal two-state (folded/unfolded) equilibrium with mean
$\so$ and variance $\so(1 - \so)$ (Bernoulli distribution). We identify this intrinsic variance with the
universal acceleration scale underlying the~\citet{Strigari2008} dSph relation:
\begin{equation}
\Gstr = \so(1 - \so)\,\ao = 0.228\,\ao = 2.74\e{-11}\;{\rm m\,s^{-2}}.
\label{eq:GStrigari}
\end{equation}
The corresponding dSph $\gc$ formula, obtained via $\gc = \gobs^2/\gbar$ at the deep-MOND limit:
\begin{equation}
\gc^{\rm dSph}(\gbar) = \frac{[\so(1 - \so)]^2 \ao^2}{\gbar}.
\label{eq:gcdSph}
\end{equation}
Of eight tested dimensionless combinations involving $\so$ and $\Tm$, only the Bernoulli variance form
matches the observed dSph-sample ensemble mean within 5\%. Alternative forms $\so\ao$ ($0.35\,\ao$,
$+47\%$ error), $(1 - \so)/\Tm$ ($0.26\,\ao$, $+10\%$), and $1/(2\Tm)$ ($0.20\,\ao$, $-15\%$) are rejected at
the $> 2\sigma$ level. Grade: B$+$ (single-population match; independent verification in Sec.~7.7).

\subsection{Tautology separation: definitional vs predictive content}
A methodological question arises: the slope $-1$ of $\log\gc$ vs $\log\gbar$ is a consequence of the
Strigari condition $\gobs \sim$ const combined with the deep-MOND extraction $\gc = \gobs^2/\gbar$, and
therefore has no independent predictive content. We separate the definitional from the predictive components
via multivariate regression on the 31-dSph sample.

\paragraph{Regression A ($\gobs$ basis).}
Regressing $\log\gobs$ on $(\log\sigma, \log\rh, \log\LV)$ with 1000-iteration bootstrap yields:
\begin{center}\small
\begin{tabular}{@{}lccll@{}}
\toprule
Coefficient & Observed & Theory & $z$-score & Verdict \\
\midrule
$\beta_\sigma$ & $+2.000 \pm 0.000$ & $+2.0$ & $+1.2$ & definitional \\
$\beta_{\rh}$ & $-1.000 \pm 0.000$ & $-1.0$ & $+0.7$ & definitional \\
$\beta_{\LV}$ & $-2\e{-7}$ & $0$ & $-1.4$ & null confirmed \\
\bottomrule
\end{tabular}
\end{center}

The fit achieves $R^2 = 1.0000$ with residual standard deviation at machine precision, recovering the Wolf
estimator $\gobs = 3\sigma^2/\rh$ as an exact identity. The $\LV$ coefficient at $2\e{-7}$ is the strongest
possible statement of $\Mbar$-independence: luminosity adds no information beyond $\sigma$ and $\rh$.

\paragraph{Orthogonalized regression.}
The dSph $\sigma$--size relation ($\log\sigma = +0.263\log\rh + $ const, consistent
with~\citealt{McConnachie2012}) introduces collinearity between $\sigma$ and $\rh$. Orthogonalizing $\sigma$
against $\rh$ and refitting preserves $\beta_{\LV} \approx 0$, providing an independent verification free of
collinearity artifacts.

\paragraph{Partial correlation.}
A third independent route: the partial correlation $r(\gobs, \LV | \sigma, \rh) = -0.236$ ($p = 0.20$) is
statistically insignificant, while $r(\gobs, \sigma | \rh, \LV)$ and $r(\gobs, \rh | \sigma, \LV)$ both
approach $\pm 1$ as required by the Wolf identity. Three independent methods (multivariate regression,
orthogonalization, partial correlation) concur: $\LV$ contributes no predictive information to $\gobs$ at
the machine-precision level.

\paragraph{Regression B ($\gc$ basis) --- quantitative confirmation of tautology.}
Fitting $\log\gc$ (where $\gc = \gobs^2/\gbar$) on the same variables yields
$(\beta_\sigma, \beta_{\rh}, \beta_{\LV}) = (+3.82, +0.16, -1.07)$, compared to the definitional expectation
$(+4, 0, -1)$ derived from $\gbar \propto \LV/\rh^2$. All three coefficients match the definitional
prediction within 0.2, confirming that the apparent slope $-1$ of $\log\gc$ vs $\log\LV$ carries zero
predictive content.

\paragraph{Consequence for the H4 hypothesis.}
The earlier H4 conjecture $\gc \propto \Mdyn^{-1}$ decomposes, via the Wolf identity, into $\gc \propto
\rh^{-3/2}$ with zero $\sigma$ contribution. The partial-correlation analysis on $\gc$ itself confirms:
$\beta_\sigma = 0.00 \pm 0.81$, $\beta_{\rh} = -1.51 \pm 0.37$. The observed ``mass scaling'' is entirely a
radius effect.

The $\rh$-dependence of $\gobs$ ($-3.4\sigma$ marginal trend) is inconsistent with strict Strigari
universality ($\gobs = $ const) and instead reflects the empirical $\sigma$--size scaling $\sigma \propto
\rh^{0.27}$~\citep{Brasseur2011, McConnachie2012}, which is not derived from the mean-field $c \to 0$
analysis. The origin of this scaling is deferred to future work.

\subsubsection{Scatter decomposition: 89\% of residual scatter is physical}
To quantify the information content of the 0.350 dex scatter in $\log(\gobs/\ao)$ around the Bernoulli
prediction, we decompose the total variance into four components: definitional (unavoidable given observed
$\sigma$ and $\rh$), measurement (Monte Carlo error propagation with 5\% $\sigma$ and 10\% $\rh$ relative
uncertainties), host environment (between-host variance across MW / M31 / Isolated populations), and
physical residual.

\begin{center}\small
\begin{tabular}{@{}lrrl@{}}
\toprule
Component & $\sigma$ (dex) & Var. fraction & Origin \\
\midrule
Total & $0.350$ & $100.0\%$ & $\log(\gobs/\ao)$ vs $\log(\Gstr/\ao)$ \\
Measurement & $0.062$ & $3.2\%$ & MC propagation, 1000 draws \\
Host (between) & $0.095$ & $7.4\%$ & MW / M31 / Isolated offsets \\
Physical residual & $0.331$ & $89.4\%$ & $\sigma$--size projection (primary) \\
\bottomrule
\end{tabular}
\end{center}

The measurement contribution is negligible (3.2\%), confirming that the scatter is not limited by
observational precision. The host-environment contribution is small (7.4\%), but its structure carries a
strong independent test:

\begin{center}\small
\begin{tabular}{@{}lrrr@{}}
\toprule
Host & $N$ & Scatter (dex) & Mean offset (dex) \\
\midrule
MW satellites & $15$ & $0.323$ & $+0.030$ \\
M31 satellites & $11$ & $0.410$ & $+0.032$ \\
Isolated & $3$ & $0.332$ & $+0.164$ \\
\bottomrule
\end{tabular}
\end{center}

The MW and M31 satellite populations yield essentially identical mean offsets from the Bernoulli prediction
($+0.030$ vs $+0.032$ dex, difference $< 5\%$ of their individual scatters). This is a non-trivial
cross-check: two independent host galaxies of comparable mass (both $\sim 10^{12}\,M_\odot$) produce
populations that calibrate $\Gstr$ to the same value, indicating that the prediction is independent of the
host galactic environment. The Isolated sample shows a marginally higher offset ($+0.164$ dex) with small
$N = 3$, in the same direction as the bridge-galaxy outer points ($+0.219\,\ao$, Sec.~7.7).

The residual physical scatter (0.331 dex) is dominated by the dSph $\sigma$--size scaling $\sigma \propto
\rh^{0.27}$, which propagates through the Wolf identity into $\gobs$. Derivation of this relation from first
principles is a candidate task for future work beyond the mean-field $c \to 0$ analysis of the present theory.

\paragraph{Summary of Section~7.6.}
(i) Three independent methods verify $\LV$-independence of $\gobs$ at machine precision, elevating the
Bernoulli prediction from population-mean agreement to a direct functional null. (ii) The $-1$ slope of
$\log\gc$ vs $\log\LV$ is quantitatively shown to be a definitional projection of the $\gbar$ construction.
(iii) The 0.350 dex residual scatter is 89\% physical and 3\% measurement; MW and M31 populations
independently calibrate $\Gstr$ to the same value, establishing host-environment independence as an
additional A-grade result of v4.7.8.

\subsection{Bridge galaxy verification: continuous C15~$\to$~Strigari transition}
For the four bridge galaxies with outer $\Ydyn > 10$ (Sec.~7.3), Rotmod\_LTG data permit radius-resolved
analysis. Radial profiles of $\gobs(R)$ and $\gbar(R)$ are constructed, and $\Ydyn(R)$ is tested for
monotonic radial increase. We define four A-grade criteria:

\begin{center}\small
\begin{tabular}{@{}lrrrrrrrr@{}}
\toprule
Galaxy & $N$ & $R$ [kpc] & $\Upsilon_{\rm in}$ & $\Upsilon_{\rm out}$ &
slope$(\log\Upsilon, \log R)$ & $z$ & MW $p$ & outer $\gobs$ [$\ao$] \\
\midrule
ESO444-G084 & 7 & 0.26--4.44 & 6.1 & 11.3 & $+0.31 \pm 0.13$ & $+2.4$ & $0.029$ & $0.317$ \\
NGC2915 & 30 & 0.34--10.04 & 11.4 & 13.2 & $+1.13 \pm 0.09$ & $+12.3$ & $0.0003$ & $0.219$ \\
NGC1705 & 14 & 0.22--6.00 & 5.4 & 12.3 & $+0.46 \pm 0.07$ & $+6.8$ & $0.0003$ & $0.282$ \\
NGC3741 & 21 & 0.23--7.00 & 5.9 & 11.5 & $+0.55 \pm 0.03$ & $+21.6$ & $0.0001$ & $0.124$ \\
\bottomrule
\end{tabular}
\end{center}

All four bridge galaxies show statistically significant radial increase of $\Ydyn$ (criterion~a: 4/4 at
$|z| > 2$) and outer-vs-inner differences confirmed by Mann--Whitney tests (criterion~b: 4/4 at $p < 0.03$).
The ensemble outer $\Ydyn$ median $= 12.76$ falls squarely in the J3 transition band (criterion~c). The
ensemble outer $\gobs$ median $= 0.219\,\ao$ agrees with the Bernoulli prediction (Eq.~\ref{eq:GStrigari},
$0.228\,\ao$) to 4\% (criterion~d). All four criteria A-grade satisfied.

\subsection{Two-population independent verification of the Bernoulli prediction}
The dSph sample (Sec.~7.5, pressure-supported systems) and the SPARC bridge galaxy outer points (Sec.~7.7,
rotation-supported outer disks) constitute two statistically independent populations. Both reproduce the
Bernoulli prediction $\Gstr = 0.228\,\ao$ at similar precision:

\begin{center}\small
\begin{tabular}{@{}lrrrl@{}}
\toprule
Population & $N$ & $\gobs$ [$\ao$] & Ratio to Bernoulli & Agreement \\
\midrule
Bernoulli prediction (Eq.~\ref{eq:GStrigari}) & --- & $0.228$ & $1.000$ & --- \\
dSph sample (Sec.~7.5) & 31 & $0.240$ & $1.053$ & 5\% \\
Bridge outer points (Sec.~7.7) & 30 & $0.219$ & $0.961$ & 4\% \\
NGC2915 outer alone & 12 & $0.219$ & $0.961$ & 4\% \\
\bottomrule
\end{tabular}
\end{center}

The two populations differ in dynamical support (pressure vs rotation), typical mass scale
($\sim 10^7\,M_\odot$ vs $\sim 10^9\,M_\odot$), typical acceleration ($\gbar \sim 10^{-3}\,\ao$ vs
$10^{-2}\,\ao$), and data reduction pathway (Wolf $\siglos$ vs Rotmod $V_{\rm bar}$). The common agreement
with a single theoretical constant derived from a two-parameter ($\Tm$, $\so$) theory, with no free fitting,
constitutes a non-trivial confirmation of the membrane $c \to 0$ limit. Grade: A (upgraded from B$+$ via
independent bridge-galaxy verification).

\subsection{Summary of dSph extension}
Results established in v4.7.8:
\begin{enumerate}[leftmargin=*,itemsep=2pt,label=(\roman*)]
\item C15 direct application to dSph fails at $+1.5$ dex (Sec.~7.2).
\item The failure is explained by J3 regime inversion $\Mjmem \ll \Mjstd$ in 28/31 dSph (A grade).
\item The $c \to 0$ limit of $U(\epsilon; c)$ yields a universal spontaneous equilibrium $\so = 0.3515$ at
$\Tm = \sqrt{6}$ (A grade).
\item The Bernoulli variance $\so(1 - \so)\ao = 0.228\,\ao$ matches the observed Strigari scale in both
dSph (5\%) and SPARC bridge outer points (4\%) with no free parameter (A grade).
\item All four bridge galaxies show significant radial transition to the J3 regime (A grade, 4/4 criteria).
\item The $\LV$-independence of $\gobs$ in dSph is verified at machine precision
($\beta_{\LV} = -2\e{-7}$, $R^2 = 1.0000$), cross-checked by orthogonalized regression and partial
correlation (A grade).
\item Scatter decomposition shows 89.4\% of the 0.350 dex residual is physical, with MW and M31 satellite
populations yielding identical mean offsets ($+0.030$ vs $+0.032$ dex) from the Bernoulli prediction,
establishing host-environment independence of $\Gstr$ (A grade).
\end{enumerate}

\textbf{Open issues:} (a) The $\rh$-dependence of $\gobs$ ($-3.4\sigma$) reflects the empirical dSph
$\sigma$--size scaling not predicted by mean-field analysis; the individual-galaxy formula
(Eq.~\ref{eq:gcdSph}) remains B-grade (0.7 dex scatter). (b) Higher-order cumulants beyond the mean-field
Bernoulli variance, or mean-field-breaking tidal dynamics, may be required to account for the residual 0.34
dex $\gobs$ scatter. (c) Ultra-faint dSph ($\Ydyn > 100$) lie at the extrapolation limit of the current
sample and warrant independent analysis.

% =============================================================================
\section{Conclusions}
% =============================================================================

We have presented a membrane cosmology framework with the following established results:

\begin{enumerate}[leftmargin=*,itemsep=4pt,label=\textbf{\arabic*.}]
\item \textbf{Geometric mean law (A-grade):} $\gc = \eta\sqrt{\ao\,G\Sigma_0}$, $\alpha = 0.5$,
$\AIC = -130$ vs MOND, 175 SPARC galaxies. Jackknife-verified (96.5\%). Origin: deep-MOND algebraic
consequence ($\alpha = 0.471 \pm 0.036$, $p(0.5) = 0.43$).

\item \textbf{MOND from first principles (A-grade):} Simple interpolation $\mu(x) = x/(1+x)$ derived from
elastic membrane potential equilibrium. Not an empirical ansatz.

\item \textbf{C15 final form (A-grade):} $\gc = 0.584\,\Yd^{-0.361}\sqrt{\ao\,\vflat^2/\hR}$.
$R^2 = 0.607$, scatter $= 0.286$ dex. MOND rejected at $p = 1.66\e{-53}$.

\item \textbf{C15 is minimal sufficient statistic (A-grade, v4.7):} Six $\Sgal$ models all rejected. All
parameter additions yield LOO $< 2\%$. 94\% of intrinsic variance irreducible. Scatter $=$ membrane fold diversity.

\item \textbf{$\kap = 0$: membrane rigidity zero (A-grade, v4.7):} $\kap \sim 10^{-40}$ kpc$^2$.
$\mathcal{L} \approx U(\epsilon; c)$. Each radius obeys $\mu(x)$ independently. Four free parameters eliminated.

\item \textbf{C14 pattern is C15 bias projection (A-grade, v4.7):} Ratio $\gc^{\rm local}/\gc^{\rm global} =
1.035$. RMS $= 0.090$ dex. C14 dynamical formulation retracted.

\item \textbf{$F_{\rm conf}$ derivation (A-grade):} $-c\ln(1 - \epsilon)$ from forest~$+$~cavity MF effective
theory (Lemmas 1--4).

\item \textbf{Weak lensing $Q > 1$ in 9/9 clusters (A-grade):} Wilcoxon $p = 0.0020$. cl1: membrane
$\AIC = -4.4$ vs NFW. Level B$+$.

\item \textbf{Deep-MOND: $\gc$ independent of $\sigbar$ (A-grade):} $r(\gc, \sigbar) = 0.07$. Membrane
responds to $\sigdyn$ ($r = 0.73$), not baryonic mass.

\item \textbf{BTFR slope is Simpson's paradox artifact (A-grade):} Within-bin slope $= 0.54 \pm 0.07$
confirms $\alpha = 0.5$ non-circularly. Dwarf $b_{gc} = 2.32$: membrane dominates rotation.

\item \textbf{Dwarf regime verified (A, v4.7.7):} SPARC-internal: 8 LITTLE THINGS dwarfs, scatter $= 0.232$
dex, KS $p = 0.67$. External: 11 non-SPARC LITTLE THINGS~\citep{Oh2015}, scatter $= 0.289$ dex, bias
$= -0.004$ dex, KS $p = 0.73$. C15 uniform across $\vflat = 12\text{--}125$ km/s. Upgraded from B$+$ via
independent confirmation.

\item \textbf{Morphological-type bridge (A-grade, v4.7):} T-type vs $\Sgal$: $\rho = +0.89$,
$p = 4\e{-60}$. Conceptual correspondence; no additional predictive power.

\item \textbf{Weak lensing RAR (A grade, v4.7.6):} KiDS-1000~\citep{Brouwer2021} constrains C15 at Mpc
scales but is CGM-limited (B grade). Independent HSC-SSP Y3 + GAMA DR4 analysis (157{,}338 isolated lenses,
503M pairs, 3 fields) yields $\gc = 2.73 \pm 0.11\,\ao$, $\AIC = +472$ favouring C15 over MOND.
$\gc$--$M_\ast$ slope ($+0.166 \pm 0.041$) rejects MOND at $4\sigma$. Slope budget: C15 ($+0.075$) + Model~B
($+0.017$) + 2-halo ($+0.034$) $= +0.126$, reducing the initial $2.2\sigma$ tension to $0.98\sigma$.

\item \textbf{Self-consistent equation and $\eta_0$ origin (A, v4.7.7):} Membrane self-consistent equation
$s = 1/(1 + \exp(-\Delta U(sQ)/\Tm))$ with $\Tm = \sqrt{6}$ yields $\langle s \rangle = 0.58$, identifying
$\eta_0 = 0.584$ as the mean elastic fraction. The 8.3\% $\beta$ discrepancy ($-0.331$ vs $-0.361$) is
resolved as finite-$\Tm$ absorption of $\Yd$ dependence by $s(Q)$, requiring no parameter adjustment.
$\Tm = \sqrt{6}$ confirmed by scatter optimisation ($\Tm^{\rm opt} = 2.35$, $0.24\sigma$ from $\sqrt{6}$).
$\Yd$ independence verified by partial correlation: $r(\gc, \Yd | \fgas) = -0.437$, $r(\gc, \fgas) = +0.06$
(n.s.).

\item \textbf{J3 regime inversion (A, v4.7.8):} Direct C15 application to 31 dSph galaxies fails with $+1.5$
dex systematic bias (four model variants consistent), not absorbed by $\Yd$ transformation. Root cause:
$\Mjmem \ll \Mjstd$ in 28/31 dSph (vs $\Mjstd \le \Mjmem$ in SPARC), inverting the causal direction of
membrane--baryon coupling. In the J3 regime, $\gc$ is set by intrinsic membrane parameters ($\Tm$, $\ao$),
not by baryonic surface density.

\item \textbf{Bernoulli prediction $\Gstr = \so(1 - \so)\,\ao$ (A, v4.7.8):} The $c \to 0$ limit of
$U(\epsilon; c)$ gives $\so = 1/(1 + \exp(3/(2\Tm))) = 0.3515$ ($\Tm = \sqrt{6}$). The thermal Bernoulli
variance predicts $\Gstr = 0.228\,\ao$ ($= 2.74\e{-11}$ m/s$^2$). Verified in two independent populations:
dSph 31-galaxy sample ($0.240\,\ao$, 5\% agreement) and bridge galaxy outer points ($0.219\,\ao$, 4\%
agreement). Non-analytic $\sqrt{Q}$ correction $s(Q) \approx \so + 0.110\sqrt{Q}$ arises from the
$\ln(1 - \epsilon)$ term.

\item \textbf{Continuous C15~$\to$~Strigari transition (A, v4.7.8):} Four SPARC bridge galaxies
(ESO444-G084, NGC2915, NGC1705, NGC3741) with outer $\Ydyn > 10$ exhibit statistically significant radial
increase of $\Ydyn(R)$ (4/4 galaxies, $z = 2.4$ to $21.6$). Mann--Whitney tests confirm outer $>$ inner
$\Upsilon$ in 4/4 galaxies ($p < 0.03$). Bridge outer median $\Ydyn = 12.76$ sits at the center of the J3
transition band ($10\text{--}30$). All four A-grade criteria satisfied: the membrane regime transitions
continuously within individual galaxies.

\item \textbf{$\gobs$ $\LV$-independence at machine precision (A, v4.7.8):} Multivariate regression
$\log\gobs$ vs $(\log\sigma, \log\rh, \log\LV)$ in 31 dSph galaxies yields $\beta_{\LV} = -2\e{-7}$ with
$R^2 = 1.0000$, recovering $\gobs = 3\sigma^2/\rh$ as an exact identity. The $\LV$ contribution is at
machine-precision level. Cross-verified by orthogonalized regression and partial correlation
$r(\gobs, \LV | \sigma, \rh) = -0.236$ ($p = 0.20$). The apparent slope $-1$ of $\log\gc$ vs $\log\LV$ is
shown to be the definitional consequence of $\gc = \gobs^2/\gbar$: regression yields
$(\beta_\sigma, \beta_{\rh}, \beta_{\LV}) = (+3.82, +0.16, -1.07)$, all within 0.2 of the definitional
prediction $(+4, 0, -1)$.

\item \textbf{Scatter decomposition and host-environment independence (A, v4.7.8):} Four-component
decomposition of the 0.350 dex $\gobs$ scatter around the Bernoulli prediction: measurement 0.062 dex
(3.2\%), host-between 0.095 dex (7.4\%), physical residual 0.331 dex (89.4\%). MW satellites ($N = 15$) and
M31 satellites ($N = 11$) yield essentially identical mean offsets from $\Gstr$ ($+0.030$ vs $+0.032$ dex),
establishing that $\Gstr$ is independent of the host galactic environment. Measurement contribution is
negligible; the physical residual is dominated by the dSph $\sigma$--size relation $\sigma \propto
\rh^{0.27}$.
\end{enumerate}

\paragraph{Retracted/revised in v4.7.8:}
Condition~14 dynamical formulation $g_{c,\rm eff}(r)$ (v4.7); $\kap \to A, B, \tau$ inverse estimation;
Euler--Lagrange derivation pathway; gradient term $\kap(d\epsilon/dr)^2$; $\Sgal$ predictive significance;
H3 hysteresis-only hypothesis for dSph offset (rejected in v4.7.8 by three independent tests); strict
Strigari universality ($\gobs = $ const) --- the $\rh$ dependence is accommodated by the empirical
$\sigma$--size relation, with $\gobs$ $\Mbar$-independence retained as the membrane prediction.

\paragraph{Priority future work:}
(1) BIG-SPARC external verification ($\sim 4000$ galaxies).
(2) Non-SPARC LITTLE THINGS galaxies: full 3.6~$\mu$m photometry for precise $\Yd$.
(3) Ultra-faint dSph verification ($\Ydyn > 100$ regime, $N > 50$ sample).
(4) Rubin LSST / Euclid weak lensing (2-halo detection at $2\text{--}3\sigma$).
(5) Independent $\Tm$ verification from non-SPARC datasets.
(6) Theoretical origin of the dSph $\sigma$--size relation $\sigma \propto \rh^{0.27}$: identify whether
higher cumulants beyond the mean-field $\so(1 - \so)$, or mean-field-breaking dynamical effects (tidal
history, formation epoch), dominate the 0.34 dex $\gobs$ scatter.

% =============================================================================
\section*{Acknowledgments}
% =============================================================================

This research is based on observations made with ESO Telescopes at the La Silla Paranal Observatory under
programme IDs 177.A-3016, 177.A-3017, 177.A-3018, and 179.A-2004, and on data products produced by the KiDS
consortium. The KiDS production team acknowledges support from: Deutsche Forschungsgemeinschaft; ERC; NOVA
and NWO-M grants; Target; the University of Padova; and the University Federico~II (Naples).

The Hyper Suprime-Cam (HSC) collaboration includes the astronomical communities of Japan and Taiwan, and
Princeton University. The HSC instrumentation and software were developed by the National Astronomical
Observatory of Japan (NAOJ), the Kavli Institute for the Physics and Mathematics of the Universe (Kavli
IPMU), the University of Tokyo, KEK, ASIAA, and Princeton University. Funding was contributed by the FIRST
program from the Japanese Cabinet Office, MEXT, JSPS, JST, the Toray Science Foundation, NAOJ, Kavli IPMU,
KEK, ASIAA, and Princeton University. GAMA is a joint European--Australasian project based around a
spectroscopic campaign using the Anglo-Australian Telescope, funded by STFC (UK), ARC (Australia), AAO, and
participating institutions.

% =============================================================================
\bibliographystyle{plainnat}
\nocite{*}  % show all references in refs_v478.bib (matches original v4.7.8 References list)
\bibliography{refs_v478}

\end{document}
